%%%%%%%%%%  scan idx 59  |  folio 48 (continued)  %%%%%%%%%%

\underline{Definición 1.19.} (1) Un objeto
$\xi\in\widehat{\underline A}$ se llama \underline{objeto formal
versal} si todo morfismo sobreyectivo $\alpha:a'\longrightarrow a$ en
$\underline A$ induce un morfismo sobreyectivo
\[
 \operatorname{Hom}_{\widehat{\underline A}}(\xi,a')
 \longrightarrow
 \operatorname{Hom}_{\widehat{\underline A}}(\xi,a).
\]

(2) Un objeto formal versal $\xi\in\widehat{\underline A}$, con
imagen $(R,s)$ en $\widehat C_{\Lambda,K}$, se llama
\underline{mínimo} si se cumple la condición siguiente. Para dos
morfismos cualesquiera
\[
 \varphi_1,\varphi_2:R\rightrightarrows D_{K'}(K),
\]
si $\varphi_1(\xi)$ es isomorfo a $\varphi_2(\xi)$, existe una
derivación $D:K\longrightarrow K'$ tal que
\[
 (1,D)(\varphi_1)=\varphi_2,
 \qquad\text{es decir,}\qquad
 \varphi_2-\varphi_1=Ds.
\]

\underline{Teorema 1.20. Si $\underline A$ es semihomogéneo y}
\[
 \dim_{K'}[\underline A(D_{K'}(K))]<\infty,
\]
\underline{entonces $\underline A$ admite un objeto formal mínimamente
versal, y dos objetos cualesquiera de esta clase son isomorfos.}

%%%%%%%%%%  scan idx 60  |  folio 49  %%%%%%%%%%

Con la notación de 1.18, sea $H^*$ un subespacio de
$[\underline A(D_{K'}(K))]$ complementario de la imagen de $\gamma$,
y sea $H$ el dual de $H^*$. Pongamos
\[
 R_0=D_H(K).
\]
Elijamos $\xi_0\in\underline A(R_0)$ cuya clase, bajo las
identificaciones
\[
 [\underline A(R_0)]
 \simeq H\otimes[\underline A(D_{K'}(K))]
 \simeq
 \operatorname{Hom}\bigl(H^*,[\underline A(D_{K'}(K))]\bigr),
\]
es la inclusión
$H^*\hookrightarrow[\underline A(D_{K'}(K))]$. Entonces:

(1) para todo $V$, los morfismos
\[
 \operatorname{Hom}_{C_{\Lambda,K}}
   \bigl(D_H(K),D_V(K)\bigr)
 \longrightarrow[\underline A(D_V(K))]
\]
son sobreyectivos;

(2) se cumple la condición de minimalidad de 1.19(2) para
$(R_0,\xi_0)$.

Elijamos $S\in\widehat C_{\Lambda,K}$ formalmente liso sobre
$\Lambda$, junto con un epimorfismo $t:S\longrightarrow R_0$, de
modo que
\[
 \Omega_{R_0/\Lambda}\otimes K'
 \xrightarrow{\ \sim\ }
 \Omega_{S/\Lambda}\otimes K'.
\]
Si $M$ es el ideal maximal de $S$, definamos por inducción, para
$n\geq1$,
\[
 R_n=S/I_n,
 \qquad
 \xi_n\in\operatorname{Ob}(\underline A(R_n)),
\]
mediante las condiciones siguientes:

(a) $I_{n+1}$ es la intersección de los ideales $I$ tales que
\[
 I_nM\subset I\subset I_n
\]
y $\xi_n$ se extiende a un objeto $\xi\in\underline A(S/I)$;

(b) $\xi_{n+1}$ es una extensión de $\xi_n$.

Si
\[
 R=\varprojlim_nR_n,
 \qquad
 \xi=\varprojlim_n\xi_n,
\]
entonces $(R,\xi)$ es un objeto formal mínimamente versal y es único
salvo isomorfismo.

